\documentclass[12pt,a4paper]{article}

% Encoding and fonts
\usepackage[T1]{fontenc}
\usepackage[utf8]{inputenc}
\usepackage{microtype}

% Page layout
\usepackage[margin=2.5cm]{geometry}
\usepackage{setspace}
\onehalfspacing

% Math
\usepackage{amsmath,amssymb}

% Tables
\usepackage{longtable,booktabs,array}
\usepackage{calc}
\usepackage{etoolbox}
\makeatletter
\patchcmd\longtable{\par}{\if@noskipsec\mbox{}\fi\par}{}{}
\makeatother
\usepackage{footnote}
\makesavenoteenv{longtable}

% Graphics
\usepackage{graphicx}
\usepackage{float}
\usepackage{placeins}
\graphicspath{{../figures/}{figures/}{./}}
\def\fps@figure{htbp}

% Links
\usepackage[hidelinks]{hyperref}
\usepackage{url}
\urlstyle{same}

% Lists
\providecommand{\tightlist}{%
  \setlength{\itemsep}{0pt}\setlength{\parskip}{0pt}}

% Paragraph formatting
\setlength{\parindent}{0pt}
\setlength{\parskip}{6pt plus 2pt minus 1pt}
\setlength{\emergencystretch}{3em}

% Section numbering
\setcounter{secnumdepth}{3}

%% ============================================================
%% Title and Authors
%% ============================================================
\title{Spatiotemporal Traffic Congestion Patterns and Network Centrality in Indonesian Metropolitan Cities}

\author{%
  Firman Hadi\textsuperscript{1,*} \and
  Yasser Wahyuddin\textsuperscript{1} \and
  L.M. Sabri\textsuperscript{1} \and
  Agung Indrajit\textsuperscript{2}%
}

\date{}

\begin{document}

\maketitle

\noindent\textsuperscript{1} Department of Geodetic Engineering, Universitas Diponegoro, Semarang, Indonesia\\
\textsuperscript{2} Deputy for Green and Digital Transformation, Nusantara Capital Authority, Indonesia\\[6pt]
\textsuperscript{*} Corresponding Author: Firman Hadi (\href{mailto:firmanhadi21@lecturer.undip.ac.id}{firmanhadi21@lecturer.undip.ac.id})

\bigskip

%% ============================================================
%% Abstract
%% ============================================================
\begin{abstract}
Urban traffic congestion presents major challenges to sustainable growth in rapidly growing Southeast Asian cities. This study offers a detailed spatiotemporal analysis of traffic congestion in three key Indonesian metropolitan areas—Jakarta, Bandung, and Semarang—using high-resolution traffic flow data from the HERE Traffic API over 11 months (March 2025 to February 2026), totaling over 264 million observations across 18,694 road segments. In addition to the empirical analysis, this study introduces a modular, open-source Python package that delivers a fully programmable, end-to-end traffic analysis pipeline—from HERE API data collection to geostatistical analysis, HPC-enabled network centrality calculations, and automated manuscript creation—eliminating dependencies on R or manual processing steps. The methodology integrates jam factor analysis with geostatistical methods (Moran's I, LISA, Getis-Ord Gi*), OpenStreetMap network analysis using OSMnx, and a novel graph-based capacity drop detection The methodology combines jam factor analysis with geostatistical tools such as Moran's I, LISA, and Getis-Ord Gi*, along with OpenStreetMap network analysis using OSMnx, and introduces a new graph-based capacity drop detection method framework. Our central finding is that \textbf{temporal factors dominate spatial factors by 1,000--4,000x}: time-of-day ($\eta^2$ = 15--24\%) vastly outweighs network centrality, POI density, and road capacity ($R^2$ < 0.01\%) in explaining relative congestion variance. Because the HERE jam factor normalizes speed to each segment's free-flow baseline, these findings characterize demand-to-capacity ratio rather than absolute delay. Triangulated evidence from four independent spatial predictors (centrality, POI density, road capacity, capacity drops) all yield null results. Road capacity shows negligible correlation with congestion (r < 0.04); spatial capacity drops show no proximity effect (p = 0.10--0.83, $|d|$ < 0.09); local bottleneck gradients are indistinguishable from surrounding segments ($|d|$ < 0.05). A spatial scale-sensitivity check using Uber H3 hexagonal aggregation at neighbourhood ($\approx$461\,m) and block ($\approx$174\,m) scales confirms that null POI and centrality findings are not a modifiable areal unit artefact; Jakarta alone exhibits weak but consistent neighbourhood-scale spatial autocorrelation (Moran's $I \approx 0.03$, $p < 0.05$) not detectable at the road-segment level. This demonstrates that congestion is fundamentally a \textbf{temporal synchronization problem}---millions of people traveling at the same times---rather than a spatial infrastructure constraint. The evening peak (16:00--19:00) shows congestion approximately 40\% above daily averages. These findings support prioritizing demand management strategies alongside infrastructure expansion, contributing empirical evidence to traffic policy debates in rapidly urbanizing contexts.
\end{abstract}

\noindent\textbf{Keywords:} Urban traffic congestion; Spatiotemporal analysis; Network centrality; Jam factor; Capacity bottleneck; Indonesian cities; HERE Traffic API; OSMnx; Spatial autocorrelation

\newpage

%% ============================================================
%% Main Body
%% ============================================================
\section{Introduction}

Urban traffic congestion poses a major challenge for rapidly growing cities in developing countries, characterized by high variability, peak periods, and reliance on informal transportation systems (Gakenheimer, 1999; Pojani \& Stead, 2015). The economic impact is substantial: in Southeast Asia, congestion costs are estimated at 2–5\% of GDP each year (Asian Development Bank, 2019), while welfare losses from transport-related air pollution amount to 4.8\% of GDP in East Asia and 3.5\% in South Asia (World Bank, 2022). The Western Pacific and Southeast Asia regions face the highest levels of traffic-related air pollution, reflecting the rapid urbanization in these developing nations (WHO, 2021). Accelerating urbanization and population growth have increased the demand for urban infrastructure, leading to persistent congestion that traditional planning methods often cannot fix (Cervero, 2013). Cities across Southeast Asia, South Asia, and Sub-Saharan Africa face similar problems: vehicle numbers grow faster than road capacity, motorcycle use is high, public transit systems are inadequate, and data for evidence-based planning is limited (Pojani \& Stead, 2015; Kutzbach, 2009). Understanding the spatial and temporal patterns of congestion in these areas is crucial for creating effective strategies, but a lack of data often limits traffic improvement efforts to places where congestion is easiest to measure.

Indonesia vividly illustrates these challenges. Jakarta has long struggled with severe traffic jams, and recent studies show that Bandung, Semarang, Surabaya, and Yogyakarta now face similar congestion issues. The World Bank reports that Jakarta is among the most traffic-congested cities worldwide, with commuters spending an average of 62 hours annually in traffic and economic losses projected to reach US\$6.5 billion by 2020. Peak-hour vehicle speeds are only about 10 km/h (World Bank, 2021). Castrol notes that Jakarta drivers experience roughly 33,240 stop-starts each year, with 27\% of travel time spent stationary (Castrol, 2015). Bandung, Indonesia’s second-largest city, ranks second among 38 Indonesian cities in a World Bank analysis, with peak-hour speeds of only 17 km/h (World Bank, 2019). The Asian Development Bank identified Bandung as Indonesia’s most congested city and 14th in Asia (ADB, 2019). A notable feature of Indonesian urban traffic is the very high motorcycle share; in Bandung, motorcycles accounted for 72.2\% of vehicles in 2013 (BPS Bandung, 2015), creating traffic flows that differ significantly from those in car-dominated systems in developed nations. Earlier Indonesian traffic research mainly relied on surveys: Susilo et al. (2007) examined motorization and public transport trends in Jakarta, while Joewono \& Kubota (2008) studied paratransit user satisfaction. More recent studies focus on demand management measures such as car-free days and odd-even plate policies in Bandung (Syaiful \& Fadlilah, 2020). However, these often provide only snapshot views rather than continuous traffic dynamics. Although Indonesia has prioritized mass transit development in six metropolitan areas through projects such as the Indonesia Mass Transit Project (World Bank, 2022), comprehensive empirical analysis of congestion patterns using high-resolution data remains limited, hindering evidence-based planning of interventions.


Recent international traffic research increasingly uses multi-source data and advanced analytical frameworks to understand urban congestion patterns. In China’s largest cities, real-time traffic data from online maps has been used to identify frequently congested road segments, showing that 1.8–9.87\% of urban road networks face chronic congestion and that congestion severity correlates with city size (Cui et al., 2024). Similar methods in Shanghai during and after the COVID-19 pandemic have shown how disruptions changed spatiotemporal congestion trends (Zhang et al., 2022). In South Asia, Dhaka has been a case study for analyzing the effects of urbanization on traffic, with researchers combining satellite-based traffic and pollution data with socioeconomic information. In Tanzania, Dar es Salaam has been examined using European Space Agency Sentinel-5P satellite data to associate traffic congestion with air pollution exposure (Dasgupta et al., 2021). These studies show that remote sensing and crowdsourced data can support detailed analyses in cities without traditional traffic-monitoring systems. Combining multi-source data with high-performance computing allows for measuring key spatiotemporal patterns in human mobility and congestion (Cai et al., 2024), although machine learning methods still face challenges in effectively merging temporal features with network topology for congestion prediction (Song et al., 2019; Ermagun \& Levinson, 2018).

The availability of probe vehicle data and APIs has advanced traffic analysis by enabling high-resolution, broad coverage (Leduc, 2008; Jenelius \& Koutsopoulos, 2015). Floating car data (FCD) from GPS vehicles offers a cost-effective alternative to fixed sensors for large-scale traffic assessment (Yuan et al., 2025). Studies show that equipping 3–8\% of vehicles with FCD and appropriate sampling yields accurate speed estimates with errors of 15 meters or less (Ma et al., 2024). Commercial APIs from HERE, TomTom, and Google are vital for urban traffic research. HERE’s data combines information from various sources across 70+ countries and 13 million km of roads (HERE Technologies, 2024). The jam factor, which ranges from 0 (free) to 10 (standstill), supports cross-network comparisons and has been used to analyze the spread of congestion in European cities (Rempe et al., 2016). Crowdsourced data research has expanded from Western cities to developing regions, such as Al Ain (UAE), building on methods from Windsor, Toronto, and New York (Alkaabi et al., 2024). However, most studies analyze a single city using snapshot data, rely on descriptive statistics, and focus more on visualization than on hypothesis testing. No studies have systematically tested whether spatial factors predict congestion while controlling for time, or compared congestion across multiple cities with hundreds of millions of observations. Research is also limited in Southeast Asia and Indonesia, where motorcycles create unique traffic dynamics.

Alongside advancements in data, Boeing’s (2017) OSMnx framework has improved capabilities for large-scale street network analysis, allowing systematic exploration of urban infrastructure topology and the calculation of transportation-related metrics, including intersection density, circuity, average node degree, and centrality measures (Boeing, 2025). Betweenness centrality—measuring the proportion of shortest paths passing through each node or edge—has been hypothesized to relate to traffic volumes, with studies finding significant links in some urban contexts (Gao et al., 2013; Porta et al., 2006). However, empirical evidence varies: research in Portland showed that grid-like networks distribute importance more evenly (15\% of paths through the most central node) than constrained networks (43\%), suggesting that topology–congestion relationships depend on network structure (Boeing, 2018). Spatial autocorrelation methods offer powerful tools for identifying traffic clusters: Global Moran’s I evaluates overall spatial clustering (Moran, 1950), while Local Indicators of Spatial Association (LISA) identify specific hotspot and coldspot locations (Anselin, 1995). Getis-Ord Gi* statistics complement LISA by detecting statistically significant spatial concentrations (Getis \& Ord, 1992; Ord \& Getis, 1995). These techniques have been widely used in traffic crash analysis—Gedamu et al. (2024) compared crash severity clustering in Addis Ababa and Berlin using Moran’s I—but applications to real-time congestion data are still less common, especially in developing country contexts.

Despite the increasing availability of probe vehicle data and commercial traffic APIs, systematic spatiotemporal analysis of urban congestion in Indonesia using high-resolution data remains limited. Past Indonesian studies relied mainly on surveys or short-term observations, lacking continuous long-term monitoring. The combination of network topology metrics with traffic flow data is still underexplored, and whether physical capacity constraints can predict congestion patterns has not been thoroughly tested in these contexts. The three cities studied here --Jakarta, Bandung, and Semarang -- represent different types of Indonesian urban environments. Jakarta, a megacity of 10.5 million residents (metro: 34 million), consistently ranks among the world’s most congested cities (TomTom Traffic Index, 2023) and has high motorcycle volumes exceeding 60\%. Bandung (2.5 million) is Java’s second-largest city, located in a highland basin with limited network development. Semarang (1.8 million), serving as Central Java’s logistics hub, stretches across coastal lowlands and southern hills, forming distinct traffic zones.

This study addresses gaps with four methods: analyzing 11 months of traffic data totaling over 264 million observations; using geostatistical techniques like Moran’s I, LISA, and Getis-Ord Gi* to identify clusters; applying OSMnx network analysis to link topology and congestion; and testing if capacity constraints, including a new graph-based capacity drop detection, predict congestion. The research offers five key contributions: a comprehensive open-source Python software package for traffic analysis; a methodology integrating commercial traffic data with open-source network analysis; a first systematic comparison of three Indonesian cities showing null results across four spatial predictors; insights that normalized congestion metrics remove capacity effects, impacting capacity–congestion studies; and practical recommendations for traffic management in Southeast Asian cities.



\section{Literature Review}

Traffic congestion in developing countries has unique characteristics, including high variability, distinct peak times, and a strong reliance on informal transportation, setting it apart from patterns observed in developed countries (Gakenheimer, 1999). The advent of big data analytics has revolutionized how we understand urban dynamics, providing unprecedented detail and insights into the complexities of city traffic systems (Batty, 2013).

In Indonesia, for example, much of the existing research has primarily relied on survey data to assess traffic conditions (Susilo et al., 2007; Joewono \& Kubota, 2008). However, this method often falls short of capturing the continuous and dynamic nature of traffic flows. Using recent technological advances, modern studies have begun integrating various data sources and advanced analytical frameworks to enhance our understanding of urban traffic behavior and congestion patterns.

The use of probe vehicle data and commercial traffic APIs has greatly enhanced traffic analysis, allowing for more detailed assessments of traffic conditions (Leduc, 2008; Jenelius \& Koutsopoulos, 2015). A key metric in this context is the jam factor, which measures traffic conditions on a scale from 0 (free-flowing traffic) to 10 (full standstill). This metric helps compare networks and has been crucial in evaluating congestion levels in various European cities (HERE Technologies, 2023; Rempe et al., 2016).

Additionally, it has been used to analyze the spatiotemporal spread of traffic congestion, providing insights into how congestion develops over time and space (Saberi et al., 2020). Boeing’s (2017) OSMnx framework has further improved our ability to conduct large-scale street network analysis, allowing researchers to systematically study urban infrastructure. Subsequent studies using this framework have revealed important findings about the relationship between traffic volumes and betweenness centrality, a measure of a node’s importance in a network (Gao et al., 2013; Porta et al., 2006).

These findings show that the structural features of a transportation network play a key role in congestion patterns and how it spreads (Kirkley et al., 2018). The traffic bottleneck theory explains that reductions in roadway capacity can lead to congestion becoming focused at specific points, a phenomenon described by Daganzo (2007) and Arnott et al. (1993).

However, empirical evidence linking static capacity features directly to city-wide congestion remains limited. To address this gap, spatial autocorrelation methods — such as Moran’s I, Local Indicators of Spatial Association (LISA), and Getis-Ord Gi* — have become important tools for identifying traffic hotspots in urban areas (Getis \& Ord, 1992; Ord \& Getis, 1995). These methods are based on Tobler’s (1970) first law of geography, which states that spatial phenomena are not randomly distributed but are influenced by their geographic context, aiding our understanding of traffic congestion patterns (Anderson, 2009; Wang et al., 2016).


\section{Study Area and Data}

\subsection{Study Area}

The three study cities represent a hierarchy of Indonesian urban centers (Table~\ref{tab:table1studyarea}).

\textbf{Table~\ref{tab:table1studyarea}.} Study area characteristics

\begin{longtable}[]{@{}
  >{\raggedright\arraybackslash}p{(\linewidth - 8\tabcolsep) * \real{0.0822}}
  >{\raggedleft\arraybackslash}p{(\linewidth - 8\tabcolsep) * \real{0.2877}}
  >{\raggedleft\arraybackslash}p{(\linewidth - 8\tabcolsep) * \real{0.1644}}
  >{\raggedleft\arraybackslash}p{(\linewidth - 8\tabcolsep) * \real{0.2466}}
  >{\raggedright\arraybackslash}p{(\linewidth - 8\tabcolsep) * \real{0.2192}}@{}}
\toprule\noalign{}
City & Population (million) & Area (km$^2$) & Traffic Segments & Urban Typology \\
\midrule\noalign{}
\endhead
\bottomrule\noalign{}
\endlastfoot
Jakarta & ~10.5 (metro: 34) & 662 & 14,549 & Megacity, National Capital \\
Bandung &~ 2.5 & 167 & 3,069 & Large City, Regional Center \\
Semarang &~ 1.8 & 373 & 1,076 & Medium City, Provincial Capital \\
\label{tab:table1studyarea}
\end{longtable}

Jakarta covers a coastal plain with both grid and organic street layouts and has a high motorcycle volume (>60\%). Bandung is situated in a highland basin with a colonial-era city center and limited network growth. Semarang extends across coastal lowlands and southern hills, resulting in different traffic patterns between the flat northern commercial areas and the hilly southern residential zones. 

\subsection{Data Collection}

Traffic flow data was collected through the HERE Traffic API at 15-minute intervals from March 2025 to February 2026, resulting in approximately 14,100 collection cycles per city (Table~\ref{tab:table2_data}). For this study, automation was achieved with a shell-scheduled R script using the hereR package. This script queried bounding-box polygons for each city, added a GMT+7 timestamp to each record, and saved each snapshot as a GeoPackage file—producing up to 48 files per day per city without manual effort. The Python package developed as part of this work (Section 4.6) offers a native Python version of this collection process, replacing hereR with direct HERE API calls, so that the entire pipeline from data collection to final outputs operates in one language.

\textbf{Table~\ref{tab:table2_data}.} Data collection summary

\begin{longtable}[]{@{}
  >{\raggedright\arraybackslash}p{(\linewidth - 8\tabcolsep) * \real{0.0870}}
  >{\raggedright\arraybackslash}p{(\linewidth - 8\tabcolsep) * \real{0.2609}}
  >{\raggedleft\arraybackslash}p{(\linewidth - 8\tabcolsep) * \real{0.1884}}
  >{\raggedleft\arraybackslash}p{(\linewidth - 8\tabcolsep) * \real{0.2174}}
  >{\raggedleft\arraybackslash}p{(\linewidth - 8\tabcolsep) * \real{0.2464}}@{}}
\toprule\noalign{}
City & Collection Period & Total Files & Total Records & Unique Segments \\
\midrule\noalign{}
\endhead
\bottomrule\noalign{}
\endlastfoot
Jakarta & ~Mar 2025 - Feb 2026 & 14,132 & 206,316,468 & 14,549 \\
Bandung & ~Mar 2025 - Feb 2026 & 14,136 & 43,407,684 & 3,069 \\
Semarang & ~Mar 2025 - Feb 2026 & 14,122 & 15,212,072 & 1,076 \\
\textbf{Total} & & \textbf{42,390} & \textbf{264,936,224} & \textbf{18,694} \\
\label{tab:table2_data}
\end{longtable}

Each record contains: road segment geometry (WGS84, EPSG:4326), jam factor (0--10), current speed, free-flow speed, confidence level, and timestamp (UTC, converted to GMT+7).

\subsection{Data Processing and Aggregation}

Raw traffic data were grouped into eight time periods that reflect Indonesian urban activity patterns (Table~\ref{tab:table3_temporal}).

\textbf{Table~\ref{tab:table3_temporal}.} Temporal period definitions

\begin{longtable}[]{@{}lll@{}}
\toprule\noalign{}
Period & Time Range & Rationale \\
\midrule\noalign{}
\endhead
\bottomrule\noalign{}
\endlastfoot
Night & 00:00-05:59 & Minimal activity \\
Morning Peak & 06:00-08:59 & Commute to work/school \\
Morning Off-Peak & 09:00-11:59 & Business hours \\
Lunch Hours & 12:00-13:59 & Midday break period \\
Afternoon Off-Peak & 14:00-16:59 & Business hours \\
Evening Peak & 16:00-18:59 & Return commute \\
Evening Off-Peak & 19:00-21:59 & Evening activities \\
Late Night & 22:00-23:59 & Reduced activity \\
\label{tab:table3_temporal}
\end{longtable}

\subsection{Network Data and Data Quality}

Street network data were collected from OpenStreetMap using OSMnx (Boeing, 2017), filtered for drivable roads (Table~\ref{tab:table4_network}). Extensive exploratory analysis verified data integrity across all 24 dataset combinations  (3 cities $\times$ 8 periods): no null values in key variables, all jam factor means between 0 and 10, and consistent observation density ($\sim$1,766 per segment). Detailed validation tables are provided in Supplementary Material S1.

\textbf{Table~\ref{tab:table4_network}.} Network statistics

\begin{longtable}[]{@{}lllll@{}}
\toprule\noalign{}
City & Nodes & Edges & Street Density (km/km$^2$) & Mean Degree \\
\midrule\noalign{}
\endhead
\bottomrule\noalign{}
\endlastfoot
Jakarta & $\sim$45,000 & $\sim$95,000 & 18.4 & 2.89 \\
Bandung & $\sim$18,000 & $\sim$38,000 & 22.1 & 2.76 \\
Semarang & $\sim$12,000 & $\sim$25,000 & 15.2 & 2.81 \\
\label{tab:table4_network}
\end{longtable}



\FloatBarrier
\section{Methodology}

\subsection{Analytical Framework}

Our framework integrates five complementary approaches (Figure~\ref{fig:fig1_analytic}): (1) descriptive statistical analysis of temporal congestion patterns; (2) geostatistical analysis of spatial clustering; (3) network analysis of topology--congestion relationships; (4) bottleneck analysis testing capacity constraint effects; and (5) spatial scale-sensitivity analysis using Uber H3 hexagonal aggregation to evaluate robustness of spatial findings to the modifiable areal unit problem.

\begin{figure}[H]
\centering
\includegraphics[width=\textwidth]{figures/analytical_framework.png}
\caption{Analytical framework integrating temporal analysis, geostatistical methods, and network topology assessment}
\label{fig:fig1_analytic}
\end{figure}

\subsection{Temporal Pattern Analysis}

Temporal patterns were examined by calculating summary statistics for each period across all road segments. One-way ANOVA tested for significant differences between periods, with post-hoc Tukey HSD tests identifying specific significant pairs.

\subsection{Geostatistical Analysis}

\subsubsection{Global Spatial Autocorrelation}

Spatial autocorrelation was assessed using Moran's I (Moran, 1950):

\[I = \frac{n}{\sum_i \sum_j w_{ij}} \cdot \frac{\sum_i \sum_j w_{ij}(x_i - \bar{x})(x_j - \bar{x})}{\sum_i (x_i - \bar{x})^2}\]

where $n$ is the number of spatial units, $w_{ij}$ is the spatial weight between units $i$ and $j$, and $x_i$ is the jam factor at unit $i$. Spatial weights used K-nearest neighbors (k=8) based on segment centroids, preferred over queen contiguity for linear road segment geometries.

\subsubsection{Local Indicators of Spatial Association}

Local Moran's I (Anselin, 1995) identified hotspot and coldspot clusters:

\[I_i = \frac{(x_i - \bar{x})}{\sigma^2} \sum_j w_{ij}(x_j - \bar{x})\]

LISA was computed with KNN (k=8) weights, 999 random permutations, and Benjamini-Hochberg FDR correction. Segments were classified as HH (hotspot), LL (coldspot), or spatial outliers (HL/LH). Temporal stability was assessed via coefficient of variation: $CV = (\sigma/\mu) \times 100\%$.

\subsection{Spatial Scale Sensitivity Analysis (H3 Hexagonal Aggregation)}
\label{sec:method_h3}

A persistent methodological concern in urban spatial analysis is the modifiable areal unit problem (MAUP; Openshaw, 1984): statistical relationships can change substantially depending on the spatial unit chosen for analysis. Road segments vary in length from under 200\,m to over 3\,km, and predictors such as POI density and network centrality may operate at a neighbourhood rather than a segment scale. To test whether null correlations at the road-segment level reflect genuine absence of relationships or a unit-of-analysis artefact, we aggregated all variables into regular hexagonal bins using the Uber H3 hierarchical geospatial indexing system (Brodsky, 2018) and re-ran the full correlation and spatial autocorrelation suite.

\textbf{Choice of hexagonal grid.} Hexagons are preferred over square grids for spatial aggregation because all six neighbours share an equal centroid-to-centroid distance, minimising directional bias in spatial weight matrices (Sahr et al., 2003). H3's hierarchical structure allows analysis at multiple scales using a consistent topology. Two resolutions were selected to span the neighbourhood and block scales relevant to urban traffic dynamics:

\begin{itemize}
  \item \textbf{Resolution 8}: mean hex diameter $\approx$461\,m, area $\approx$0.46\,km$^2$ (neighbourhood scale; 235--1,836 hexagons per city)
  \item \textbf{Resolution 9}: mean hex diameter $\approx$174\,m, area $\approx$0.10\,km$^2$ (block scale; 490--5,319 hexagons per city)
\end{itemize}

\textbf{Variable aggregation.} Each road segment centroid was assigned to its containing H3 cell using \texttt{h3.latlng\_to\_cell()}. Jam factor per hexagon was computed as the observation-weighted mean across all segments and all eight temporal periods within the cell:

\[\overline{JF}_{\text{hex}} = \frac{\displaystyle\sum_{s \in \text{hex}} \sum_{t} JF_{s,t} \cdot n_{s,t}}{\displaystyle\sum_{s \in \text{hex}} \sum_{t} n_{s,t}}\]

where $s$ indexes road segments, $t$ indexes temporal periods, $JF_{s,t}$ is the mean jam factor for segment $s$ in period $t$, and $n_{s,t}$ is the corresponding observation count. POI count per hexagon was the number of OSM POI centroids falling within each cell. Mean betweenness centrality per hexagon was the arithmetic mean of all OSM edge centrality values whose midpoints fell within the cell.

\textbf{Statistical tests.} Pearson and Spearman correlations were computed between hex-level jam factor and each predictor variable. Global Moran's I was computed using Queen contiguity weights derived from hexagon polygon adjacency---geometrically well-defined for regular hexagons, where each interior cell has exactly six neighbours---with row-standardisation and 999 random permutations for inference. Results are reported alongside the corresponding segment-level baseline values to quantify the effect of spatial scale on each statistic.

\subsection{Network Centrality Analysis}

Following Boeing (2017), we computed two network centrality metrics. \textbf{Betweenness centrality} measures the fraction of shortest paths passing through each edge:

\[C_B(e) = \sum_{s \neq t} \frac{\sigma_{st}(e)}{\sigma_{st}}\]

\textbf{Closeness centrality} measures average distance to all other nodes:

\[C_C(v) = \frac{n-1}{\sum_{u \neq v} d(u,v)}\]

Edge betweenness was computed on undirected networks weighted by edge length, with approximate sampling (k=500 source nodes) for large networks.

\subsection{Bottleneck and Capacity Drop Analysis}

To determine if physical capacity constraints cause congestion, we created a multi-scale framework. Road segments received ordinal capacity scores (1–6) based on OSM highway classifications: motorway (6) through tertiary (2), with link variants scoring half their parent. Networks were filtered to include only HERE-comparable road types before analysis.

\subsubsection{Graph-Based Capacity Drop Detection}

Capacity drops—where the road hierarchy decreases along the travel direction—were identified using the directed network graph. For each node, we compared the maximum capacity scores of incoming and outgoing edges.

\[\text{Drop ratio} = \frac{C_{\text{in}} - C_{\text{out}}}{C_{\text{in}}}\]

Nodes with a drop ratio of  $\geq$ 0.20 were classified as capacity drop locations. The Euclidean distance from each traffic segment to the nearest drop was calculated.

\subsubsection{Local Bottleneck Gradient and Statistical Testing}

A local bottleneck gradient used K=10 nearest network neighbors per segment: segments with lower capacity than at least one neighbor were classified as “local bottleneck,” while others were classified as “non-bottleneck.” Three complementary tests assessed the capacity–congestion relationship: (1) a two-sample t-test comparing high- versus low-capacity road segments (Cohen’s d); (2) Pearson and Spearman correlations between distance-to-nearest-drop and jam factor, plus a median-split t-test; and (3) a t-test comparing local bottleneck versus non-bottleneck segments.

\subsection{Python Package and Computational Pipeline}

A key contribution of this study is a modular, open-source Python package that offers a  a \textbf{fully programmable, end-to-end traffic analysis pipeline}---covering data collection, aggregation, spatial analysis, and manuscript creation entirely within Python, without depending on R or other languages. The package includes about 6,700 lines of code organized into four stages across 20 scripts:

\textbf{Stage 1 -- Python-native data collection.} The package performs direct HERE Traffic API calls in Python, replacing the R/hereR method used to gather data for this study. The collector queries configurable bounding-box polygons at user-defined intervals (default: 15 minutes), adds a timezone-aware GMT+7 timestamp, and saves each snapshot as a GeoPackage file—producing up to 96 files per city daily with no manual effort. City configurations (bounding boxes, expected segment counts, output paths) are centralized in  \texttt{config.py}, making the collector easily extendable to any city supported by the HERE API.

\textbf{Stage 2 -- Large-scale aggregation.} City-specific aggregation scripts process all raw GeoPackage files, match road segments using MD5-hashed geometries implemented in \texttt{utils.py}, and compute per-segment summary statistics (mean, std, min, max, count) for each of the eight temporal periods. During the 11-month study, this stage handled over 42,000 input files and 264 million traffic records.

\textbf{Stage 3 -- Analysis.} Ten dedicated scripts perform the complete analytical workflow: exploratory validation  (\texttt{exploratory\_data\_analysis.py}), geostatistical analysis (\texttt{geostatistical\_analysis.py} and \texttt{geostatistical\_visualizations.py}), advanced spatial statistics including Moran's I, LISA, and spatial regression (\texttt{advanced\_spatial\_analysis.py}, \texttt{revision\_analyses.py}), bottleneck and capacity drop detection (\texttt{bottleneck\_analysis.py}), POI and activity-centre analysis (\texttt{activity\_center\_analysis.py}, \texttt{poi\_congestion\_analysis.py}), comparative temporal--spatial analysis (\texttt{temporal\_vs\_spatial\_comparison.py}), and network centrality computation (\texttt{compute\_centrality\_correlations.py}), and H3 hexagonal scale-sensitivity analysis (\texttt{h3\_robustness\_check.py}, \texttt{generate\_h3\_figures.py}). Network centrality, involving $O(VE)$ shortest-path computation over graphs of up to 95,000 edges, is offloaded via SLURM job submission (\texttt{run\_hpc.sh}) to an NVIDIA H100 GPU cluster, enabling overnight batch processing across all three cities.

\textbf{Stage 4 -- Reproducible manuscript generation.} Scripts in the \texttt{paper/} sub-package automate figure generation (\texttt{generate\_figures.py}), LaTeX compilation (\texttt{build\_latex.py}), and PDF creation (\texttt{create\_pdf.py}), so any change to analysis results propagates automatically to the manuscript. All outputs are serialised to CSV (tables) and PNG (figures), ensuring full reproducibility.

Core Python dependencies: GeoPandas, Shapely, PyProj (spatial I/O); OSMnx 2.0, NetworkX (street network analysis); esda, libpysal, spreg (spatial statistics and regression); h3 v4.4 (Uber H3 hexagonal indexing); SciPy, statsmodels (hypothesis testing); Matplotlib, Seaborn, contextily (visualisation); requests (HERE API collection). The package is available at \url{https://github.com/firmanhadi21/traffic-analyses}.



\section{Results}

\subsection{Overall Congestion Characteristics}

Table~\ref{tab:table5_congestion} presents summary statistics across all temporal periods. Jakarta exhibits the highest congestion: during the evening peak, 53.5\% of segments reach moderate congestion (JF > 2.0), compared with 11.9\% in Bandung and none in Semarang.

\textbf{Table~\ref{tab:table5_congestion}.} Congestion summary statistics by city (all temporal periods)

\begin{longtable}[]{@{}llll@{}}
\toprule\noalign{}
Statistic & Jakarta & Bandung & Semarang \\
\midrule\noalign{}
\endhead
\bottomrule\noalign{}
\endlastfoot
Mean Jam Factor & 1.43 & 1.36 & 1.19 \\
Std. Deviation & 0.47 & 0.50 & 0.40 \\
Median & 1.63 & 1.52 & 1.34 \\
Max Segment Mean & 2.25 & 2.29 & 1.98 \\
Moderate+ congestion (evening peak) & 53.5\% & 11.9\% & 0.0\% \\
\label{tab:table5_congestion}
\end{longtable}

\subsection{Temporal Patterns}

\subsubsection{Period-Based Analysis}

Figure~\ref{fig:fig2_meanjam} and Table~\ref{tab:table6_meanjam} present mean jam factors by temporal period. The evening peak (16:00--19:00) shows the highest congestion across all cities---approximately 40\% above daily averages (39--41\%). All cities follow the same diurnal progression: minimal congestion at night (JF $\approx$ 0.5), rising through the morning peak, sustained at midday, and culminating in the evening peak before declining.

\begin{figure}
\centering
\includegraphics[width=\textwidth]{figures/temporal_pattern_comparison.png}
\caption{Mean jam factor by temporal period across three cities}
\label{fig:fig2_meanjam}
\end{figure}

\textbf{Table ~\ref{tab:table6_meanjam}.} Mean jam factor by temporal period

\begin{longtable}[]{@{}llll@{}}
\toprule\noalign{}
Period & Jakarta & Bandung & Semarang \\
\midrule\noalign{}
\endhead
\bottomrule\noalign{}
\endlastfoot
Night & 0.50 & 0.46 & 0.45 \\
Morning Peak & 1.19 & 1.15 & 0.95 \\
Morning Off-Peak & 1.63 & 1.63 & 1.39 \\
Lunch Hours & 1.66 & 1.70 & 1.45 \\
Afternoon Off-Peak & 1.79 & 1.87 & 1.52 \\
Evening Peak & \textbf{2.01} & \textbf{1.92} & \textbf{1.65} \\
Evening Off-Peak & 1.67 & 1.40 & 1.33 \\
Late Night & 0.97 & 0.77 & 0.78 \\
\label{tab:table6_meanjam}
\end{longtable}

\subsubsection{Statistical Significance of Temporal Variation}

One-way ANOVA confirms that temporal period differences are highly statistically significant (Table~\ref{tab:table_anova}). Post-hoc Tukey HSD tests reveal that all 28 pairwise period comparisons are significant (p < 0.05) for each city. The effect size ($\eta^2$) ranges from 15--24\%---a strong effect for a single categorical variable. Weekday congestion exceeds weekend levels by approximately 35--40\%, with Friday evenings peaking and Sunday mornings at a minimum.

\textbf{Table~\ref{tab:table_anova}.} ANOVA results for temporal period effects

\begin{longtable}[]{@{}llll@{}}
\toprule\noalign{}
City & F-statistic & p-value & Significant Period Pairs \\
\midrule\noalign{}
\endhead
\bottomrule\noalign{}
\endlastfoot
Jakarta & 1,191,699.66 & < 0.001 & 28/28 \\
Bandung & 224,445.67 & < 0.001 & 28/28 \\
Semarang & 45,863.33 & < 0.001 & 28/28 \\
\label{tab:table_anova}
\end{longtable}

\subsection{Spatial Pattern Analysis}

\subsubsection{Global Spatial Autocorrelation}

Global Moran's I (KNN k=8 weights) is close to zero and non-significant for all cities (Table~\ref{tab:table_moran}), indicating no strong global spatial autocorrelation when aggregated across the full study period. Results are robust across alternative weight specifications (KNN k=4, k=12; distance bands 500m and 1000m), and computing Moran's I separately for each temporal period also yields non-significant results across all cities and periods (all p > 0.05). This confirms the null global autocorrelation finding is not an artifact of temporal averaging or weight specification. A spatial scale-sensitivity analysis using Uber H3 hexagonal aggregation (Section~\ref{sec:h3_robustness}) further confirms this null result for Semarang and Bandung, while revealing a weak but statistically significant neighbourhood-level autocorrelation for Jakarta ($I = 0.030$, $p = 0.033$ at resolution~8) not detectable at the road-segment level.

\textbf{Table~\ref{tab:table_moran}.} Global spatial autocorrelation statistics (Moran's I)

\begin{longtable}[]{@{}lllll@{}}
\toprule\noalign{}
City & Moran's I & Z-score & p-value & Interpretation \\
\midrule\noalign{}
\endhead
\bottomrule\noalign{}
\endlastfoot
Jakarta & 0.0026 & 0.69 & 0.492 & Random \\
Bandung & 0.0075 & 0.93 & 0.353 & Random \\
Semarang & -0.0039 & -0.21 & 0.837 & Random \\
\label{tab:table_moran}
\end{longtable}

\subsubsection{Local Hotspot Identification}

Despite weak global autocorrelation, LISA identifies meaningful local clusters (Table~\ref{tab:table_lisa}). Approximately 10\% of segments in Jakarta and Bandung participate in significant clusters (7.2\% in Semarang). After Benjamini-Hochberg FDR correction, no segments remain significant, indicating that while observed clusters exceed chance expectations by approximately 2$\times$, individual p-values are not sufficiently small to survive multiple testing correction. Uncorrected LISA results should therefore be interpreted as suggestive of local clustering tendencies; the patterns are geographically coherent (concentrated along known corridors) and exceed chance expectations by $\sim$2$\times$, suggesting genuine but weak local spatial structure. Getis-Ord Gi* analysis corroborates LISA: 9--10\% of segments are classified as significant hot or cold spots, with cold spots slightly outnumbering hot spots across all cities.

\textbf{Table~\ref{tab:table_lisa}.} LISA cluster classification (evening peak, p < 0.05)

\begin{longtable}[]{@{}
  >{\raggedright\arraybackslash}p{(\linewidth - 6\tabcolsep) * \real{0.3333}}
  >{\raggedleft\arraybackslash}p{(\linewidth - 6\tabcolsep) * \real{0.2143}}
  >{\raggedleft\arraybackslash}p{(\linewidth - 6\tabcolsep) * \real{0.2143}}
  >{\raggedleft\arraybackslash}p{(\linewidth - 6\tabcolsep) * \real{0.2381}}@{}}
\toprule\noalign{}
Cluster Type & Jakarta & Bandung & Semarang \\
\midrule\noalign{}
\endhead
\bottomrule\noalign{}
\endlastfoot
HH (Hotspot) & 438 (3.0\%) & 86 (2.8\%) & 19 (1.8\%) \\
LL (Coldspot) & 347 (2.4\%) & 73 (2.4\%) & 23 (2.1\%) \\
HL (High-Low Outlier) & 334 (2.3\%) & 69 (2.2\%) & 19 (1.8\%) \\
LH (Low-High Outlier) & 416 (2.9\%) & 92 (3.0\%) & 17 (1.6\%) \\
Not Significant & 13,014 (89.5\%) & 2,749 (89.6\%) & 998 (92.8\%) \\
\textbf{Total Significant} & \textbf{1,535 (10.5\%)} & \textbf{320 (10.4\%)} & \textbf{78 (7.2\%)} \\
\label{tab:table_lisa}
\end{longtable}

\textbf{Jakarta} hotspots concentrate in: the Sudirman-Thamrin CBD corridor; the Tangerang-Jakarta western approach; and the Bekasi-Jakarta eastern approach. \textbf{Bandung} hotspots include the Dago/Setiabudhi northern corridors, Asia-Afrika city center, and Soekarno-Hatta eastern industrial zone. \textbf{Semarang} hotspots appear along Pandanaran/Pemuda central corridors, Siliwangi coastal road, and the Majapahit eastern approach.

\begin{figure}
\centering
\includegraphics[width=\textwidth]{figures/jkt_hotspots_evening_peak.png}
\caption{Spatial distribution of congestion hotspots --- Jakarta}
\end{figure}

\begin{figure}
\centering
\includegraphics[width=\textwidth]{figures/bdg_hotspots_evening_peak.png}
\caption{Spatial distribution of congestion hotspots --- Bandung}
\end{figure}

\begin{figure}
\centering
\includegraphics[width=\textwidth]{figures/smg_hotspots_evening_peak.png}
\caption{Spatial distribution of congestion hotspots --- Semarang}
\end{figure}

\subsection{Network Topology--Congestion Relationships}

\subsubsection{Betweenness Centrality}

Traffic segments were spatially joined to OSMnx edges using centroid-based nearest-neighbor matching (200m threshold). Match quality is high: median centroid distances are under 20m, with over 85\% of matches within 50m. Correlation analysis reveals \textbf{negligible correlations} between betweenness centrality and jam factor across all cities (Table 10). Despite Bandung's statistically significant Spearman $\rho$ = 0.040 (p = 0.009), the effect size is trivially small, contradicting the common assumption that topologically important roads experience greater congestion.

\textbf{Table 10.} Centrality--congestion correlations

\begin{longtable}[]{@{}llllll@{}}
\toprule\noalign{}
City & n matched & Pearson r & p-value & Spearman $\rho$ & p-value \\
\midrule\noalign{}
\endhead
\bottomrule\noalign{}
\endlastfoot
Jakarta & 20,822 & 0.002 & 0.828 & -0.001 & 0.946 \\
Bandung & 4,336 & 0.012 & 0.442 & 0.040 & 0.009 \\
Semarang & 1,507 & -0.011 & 0.667 & -0.030 & 0.252 \\
\end{longtable}

\begin{figure}
\centering
\includegraphics[width=\textwidth]{figures/jkt_traffic_maps.png}
\caption{Edge betweenness centrality --- Jakarta}
\end{figure}

\begin{figure}
\centering
\includegraphics[width=\textwidth]{figures/bdg_traffic_maps.png}
\caption{Edge betweenness centrality --- Bandung}
\end{figure}

\begin{figure}
\centering
\includegraphics[width=\textwidth]{figures/smg_traffic_maps.png}
\caption{Edge betweenness centrality --- Semarang}
\end{figure}

\subsubsection{POI Density Analysis}

Point of Interest density was computed using network-distance analysis (400m walking network distance, approximately a 5-minute walk; Moudon et al., 2006). Sensitivity analyses across 200m, 400m, 800m, and 1200m thresholds confirm results are robust. POI density shows \textbf{no meaningful correlation} with congestion (Table 10a). Categorical comparison of high-POI activity centers (top quartile) versus peripheral areas (bottom quartile) confirms this: differences are negligible ($|d|$ < 0.05, p > 0.3) for all cities, demonstrating that congestion is location-independent. Spatial regression (OLS, Spatial Lag, and Spatial Error models) corroborates these results, with $R^2$ < 0.3\% across all specifications; accounting for spatial structure reveals no hidden centrality--congestion relationships.

\textbf{Table 10a.} POI density--congestion correlations

\begin{longtable}[]{@{}llllll@{}}
\toprule\noalign{}
City & Total POIs & Pearson r & p-value & Spearman $\rho$ & p-value \\
\midrule\noalign{}
\endhead
\bottomrule\noalign{}
\endlastfoot
Jakarta & --- & 0.008 & 0.284 & -0.001 & 0.946 \\
Bandung & --- & 0.005 & 0.770 & 0.004 & 0.820 \\
Semarang & --- & -0.012 & 0.640 & -0.030 & 0.252 \\
\end{longtable}

\subsection{Robustness Check: Spatial Scale Sensitivity}
\label{sec:h3_robustness}

To address the potential modifiable areal unit problem (MAUP)---the possibility that null correlations reflect a mismatch between the road-segment unit of analysis and the neighbourhood scale at which POI density and network centrality actually operate---we re-ran all spatial correlations after aggregating variables into Uber H3 hexagonal bins at two resolutions: resolution~8 (mean hex diameter $\approx$461\,m, area $\approx$0.46\,km$^2$, neighbourhood scale) and resolution~9 (diameter $\approx$174\,m, area $\approx$0.10\,km$^2$, block scale). For each hexagon, jam factor was computed as the observation-weighted mean across all road segments and temporal periods falling within the cell (235--1,836 hexagons per city at resolution~8; 490--5,319 at resolution~9). POI counts and mean edge betweenness centrality were aggregated to the same hexagonal grid using the same OSM data sources as the segment-level analyses.

POI-density and centrality correlations remained negligible at both resolutions across all three cities (all $|\rho| < 0.08$, all $p > 0.07$; Table~\ref{tab:h3_robustness}). The maximum observed effect was Semarang's POI Spearman $\rho = +0.077$ ($p = 0.242$) at resolution~8, still far below conventional significance thresholds and practically negligible in magnitude. These results confirm that the null segment-level findings reported above are not an artefact of spatial scale: POI density and network centrality do not predict congestion at any spatial grain tested, from individual road segments to 461\,m neighbourhood hexagons.

Global Moran's I on hex-aggregated jam factor was non-significant for Semarang and Bandung at both resolutions ($p > 0.05$), consistent with the segment-level finding. Jakarta, however, exhibited weak but statistically significant positive spatial autocorrelation at the neighbourhood scale at both resolutions ($I = 0.030$, $p = 0.033$ at resolution~8; $I = 0.022$, $p = 0.030$ at resolution~9), a pattern not detectable at the road-segment level ($I = 0.003$, $p = 0.492$). This suggests that Jakarta's congestion has coherent neighbourhood-level spatial structure---hexagons with elevated congestion tend to be surrounded by similarly elevated neighbours---but this clustering is not explained by POI density or betweenness centrality, both of which remain uncorrelated with congestion even within Jakarta at the hex scale. The drivers of this weak but consistent neighbourhood clustering in Jakarta warrant further investigation, with candidate explanatory variables including land-use mix, population density, and access to public transit.

\textbf{Table~\ref{tab:h3_robustness}.} H3 hex-level robustness check --- Spearman $\rho$ and Global Moran's I (p-values in parentheses)

\begin{longtable}[]{@{}llrlll@{}}
\toprule\noalign{}
City & Resolution & Hexagons & POI $\rho$ (p) & Centrality $\rho$ (p) & Moran's I (p) \\
\midrule\noalign{}
\endhead
\bottomrule\noalign{}
\endlastfoot
Semarang & 8 ($\approx$461\,m) & 235  & +0.077 (0.242) & +0.053 (0.427) & $-$0.081 (0.087) \\
Semarang & 9 ($\approx$174\,m) & 490  & +0.033 (0.462) & +0.022 (0.630) & +0.004 (0.443) \\
Bandung  & 8 ($\approx$461\,m) & 517  & +0.053 (0.228) & +0.041 (0.352) & +0.052 (0.055) \\
Bandung  & 9 ($\approx$174\,m) & 1,222 & +0.027 (0.339) & +0.053 (0.070) & $-$0.008 (0.412) \\
Jakarta  & 8 ($\approx$461\,m) & 1,836 & +0.021 (0.362) & $-$0.007 (0.784) & +0.030 (0.033)\textsuperscript{*} \\
Jakarta  & 9 ($\approx$174\,m) & 5,319 & $-$0.014 (0.314) & $-$0.003 (0.836) & +0.022 (0.030)\textsuperscript{*} \\
\label{tab:h3_robustness}
\end{longtable}

\textit{Segment-level baseline for comparison: POI $\rho < 0.013$ (all $p > 0.6$); Centrality $\rho < 0.012$ (all $p > 0.44$); Moran's I $< 0.008$ (all $p > 0.35$). \textsuperscript{*}$p < 0.05$.}

\begin{figure}[htbp]
\centering
\includegraphics[width=\textwidth]{figures/h3_congestion_maps.png}
\caption{Neighbourhood-scale congestion distribution at H3 resolution 8 ($\approx$461\,m hex diameter). Each hexagon shows the observation-weighted mean jam factor aggregated across all 8 temporal periods and all road segments within the cell. Jakarta's higher and more spatially structured congestion is visible relative to Bandung and Semarang.}
\label{fig:h3_congestion_maps}
\end{figure}

\begin{figure}[htbp]
\centering
\includegraphics[width=0.75\textwidth]{figures/h3_morans_scale.png}
\caption{Global Moran's I by spatial scale for three cities. Bars show Moran's I at the road-segment level, H3 resolution~9 ($\approx$174\,m), and H3 resolution~8 ($\approx$461\,m). Asterisk (*) denotes $p < 0.05$. Jakarta's spatial autocorrelation reaches significance only at the hexagonal aggregation scales, while Semarang and Bandung remain statistically random at all three scales.}
\label{fig:h3_morans_scale}
\end{figure}

\begin{figure}[htbp]
\centering
\includegraphics[width=\textwidth]{figures/h3_correlation_panels.png}
\caption{Scatter plots of POI density (top row) and network centrality (bottom row) versus hex-level mean jam factor at H3 resolution~8. Dashed lines show OLS regression fits. Pearson $r$ and Spearman $\rho$ are annotated per panel; no significant correlation is detected in any city for either predictor.}
\label{fig:h3_correlation_panels}
\end{figure}

\begin{figure}[htbp]
\centering
\includegraphics[width=0.85\textwidth]{figures/h3_correlation_comparison.png}
\caption{Spearman $\rho$ for POI density (left) and network centrality (right) at three spatial scales: road segment, H3 resolution~9 ($\approx$174\,m), and H3 resolution~8 ($\approx$461\,m). The grey band marks the negligible-effect region ($|\rho| < 0.1$). Correlations remain within the negligible band across all cities and scales, confirming that the null segment-level findings are not a modifiable areal unit artefact.}
\label{fig:h3_correlation_comparison}
\end{figure}

\subsection{Bottleneck and Capacity Drop Analysis}

\subsubsection{Aggregate Capacity--Congestion Relationship}

Spatial join between HERE segments and filtered OSMnx networks achieved 95--98\% match rates. Contrary to the bottleneck hypothesis, \textbf{high-capacity roads show marginally higher congestion} than low-capacity roads in all cities (Table 10c)---the opposite of capacity constraint predictions. All differences are non-significant (p > 0.05) with trivial effect sizes ($|d|$ < 0.08). This reversal reflects the HERE jam factor's normalization to free-flow speed: capacity effects are divided out by design. Pearson and Spearman correlations between capacity score and jam factor confirm near-zero relationships (Table 10d).

\textbf{Table 10c.} Congestion comparison: Low- vs.\ high-capacity road segments

\begin{longtable}[]{@{}
  >{\raggedright\arraybackslash}p{(\linewidth - 12\tabcolsep) * \real{0.0732}}
  >{\raggedright\arraybackslash}p{(\linewidth - 12\tabcolsep) * \real{0.2073}}
  >{\raggedright\arraybackslash}p{(\linewidth - 12\tabcolsep) * \real{0.2317}}
  >{\raggedright\arraybackslash}p{(\linewidth - 12\tabcolsep) * \real{0.1463}}
  >{\raggedright\arraybackslash}p{(\linewidth - 12\tabcolsep) * \real{0.0976}}
  >{\raggedright\arraybackslash}p{(\linewidth - 12\tabcolsep) * \real{0.1098}}
  >{\raggedright\arraybackslash}p{(\linewidth - 12\tabcolsep) * \real{0.1341}}@{}}
\toprule\noalign{}
City & Low Capacity JF & High Capacity JF & Difference & t-stat & p-value & Cohen's d \\
\midrule\noalign{}
\endhead
\bottomrule\noalign{}
\endlastfoot
Semarang & 1.651 & 1.655 & -0.26\% & -0.87 & 0.382 & -0.056 \\
Bandung & 1.920 & 1.925 & -0.25\% & -1.84 & 0.066 & -0.072 \\
Jakarta & 2.005 & 2.007 & -0.10\% & -1.88 & 0.060 & -0.033 \\
\end{longtable}

\textbf{Table 10d.} Capacity score--congestion correlations

\begin{longtable}[]{@{}lllll@{}}
\toprule\noalign{}
City & Pearson r & p-value & Spearman $\rho$ & p-value \\
\midrule\noalign{}
\endhead
\bottomrule\noalign{}
\endlastfoot
Semarang & 0.027 & 0.394 & 0.031 & 0.326 \\
Bandung & 0.034 & 0.074 & 0.032 & 0.087 \\
Jakarta & 0.017 & 0.051 & 0.016 & 0.066 \\
\end{longtable}

\subsubsection{Spatial Capacity Drop Analysis}

Graph-based detection identified capacity drops scaling with network size (Table 10e). Proximity to drops shows \textbf{no significant effect on congestion} (p = 0.095--0.830, $|d|$ < 0.09). Distance-to-drop correlations are similarly negligible (Pearson r: $-$0.008 to 0.019).

\textbf{Table 10e.} Capacity drop detection and proximity analysis

\begin{longtable}[]{@{}
  >{\raggedright\arraybackslash}p{(\linewidth - 12\tabcolsep) * \real{0.0789}}
  >{\raggedright\arraybackslash}p{(\linewidth - 12\tabcolsep) * \real{0.2105}}
  >{\raggedright\arraybackslash}p{(\linewidth - 12\tabcolsep) * \real{0.1711}}
  >{\raggedright\arraybackslash}p{(\linewidth - 12\tabcolsep) * \real{0.1711}}
  >{\raggedright\arraybackslash}p{(\linewidth - 12\tabcolsep) * \real{0.1053}}
  >{\raggedright\arraybackslash}p{(\linewidth - 12\tabcolsep) * \real{0.1184}}
  >{\raggedright\arraybackslash}p{(\linewidth - 12\tabcolsep) * \real{0.1447}}@{}}
\toprule\noalign{}
City & Capacity Drops & Near-Drop JF & Far-Drop JF & t-stat & p-value & Cohen's d \\
\midrule\noalign{}
\endhead
\bottomrule\noalign{}
\endlastfoot
Semarang & 78 & 1.649 & 1.656 & -1.08 & 0.283 & -0.083 \\
Bandung & 245 & 1.922 & 1.923 & -0.22 & 0.830 & -0.010 \\
Jakarta & 1,443 & 2.008 & 2.006 & 1.67 & 0.095 & 0.035 \\
\end{longtable}

\begin{figure}
\centering
\includegraphics[width=\textwidth]{figures/capacity_drop_spatial_analysis.png}
\caption{Spatial capacity drop analysis: capacity drop proximity and local bottleneck effects across three cities}
\end{figure}

\subsubsection{Local Bottleneck Gradient}

The local gradient analysis (K=10 nearest neighbors) provides the finest-grained test (Table 10f). Segments classified as local bottlenecks exhibit \textbf{no meaningful congestion difference} from non-bottleneck segments (all p > 0.4, $|d|$ < 0.05).

\textbf{Table 10f.} Local bottleneck gradient analysis

\begin{longtable}[]{@{}llllll@{}}
\toprule\noalign{}
City & Local BN JF & Non-BN JF & t-stat & p-value & Cohen's d \\
\midrule\noalign{}
\endhead
\bottomrule\noalign{}
\endlastfoot
Semarang & 1.650 & 1.654 & -0.65 & 0.517 & -0.047 \\
Bandung & 1.920 & 1.922 & -0.62 & 0.534 & -0.027 \\
Jakarta & 2.007 & 2.006 & 0.80 & 0.426 & 0.016 \\
\end{longtable}

\subsubsection{HERE Jam Factor Normalization}

A fundamental property of the jam factor shapes all spatial analyses. The metric is computed as:

\[JF \propto 1 - \frac{v_{\text{current}}}{v_{\text{free-flow}}}\]

Because each segment's congestion is measured relative to its own free-flow baseline, road class effects are divided out by design. Empirically, all road types show mean jam factors within a 4\% range (Semarang: 1.644--1.711 across all highway classes), confirming this normalization. This means the jam factor measures \textbf{demand relative to capacity} rather than absolute congestion, and any investigation of the capacity--congestion relationship using absolute terms must use unnormalized speed data.

\subsection{Temporal vs Spatial Predictors: A Critical Comparison}

Table 11a presents variance explained by each predictor type. Temporal period explains \textbf{1,000--4,000$\times$ more variance} than any spatial predictor. We note that $\eta^2$ (ANOVA) and $R^2$ (bivariate correlation) are not strictly comparable, but the magnitude of the difference (three orders of magnitude) cannot be attributed to methodological artifacts.

\textbf{Table 11a.} Variance explained by predictor type

\begin{longtable}[]{@{}
  >{\raggedright\arraybackslash}p{(\linewidth - 8\tabcolsep) * \real{0.2292}}
  >{\raggedright\arraybackslash}p{(\linewidth - 8\tabcolsep) * \real{0.1875}}
  >{\raggedright\arraybackslash}p{(\linewidth - 8\tabcolsep) * \real{0.1875}}
  >{\raggedright\arraybackslash}p{(\linewidth - 8\tabcolsep) * \real{0.1875}}
  >{\raggedright\arraybackslash}p{(\linewidth - 8\tabcolsep) * \real{0.2083}}@{}}
\toprule\noalign{}
Predictor & Measure & Jakarta & Bandung & Semarang \\
\midrule\noalign{}
\endhead
\bottomrule\noalign{}
\endlastfoot
\textbf{Time Period} & $\eta^2$ (ANOVA) & 23.8\% & 19.2\% & 15.4\% \\
POI Density & $R^2$ & 0.006\% & 0.003\% & 0.014\% \\
Network Centrality & $R^2$ & 0.0004\% & 0.014\% & 0.012\% \\
Road Capacity & $R^2$ & 0.028\% & 0.115\% & 0.072\% \\
Capacity Drop Proximity & $R^2$ & 0.006\% & 0.002\% & 0.001\% \\
\textbf{Ratio (Temporal/Spatial)} & --- & \textbf{$\sim$4000x} & \textbf{$\sim$1400x} & \textbf{$\sim$1100x} \\
\end{longtable}

\begin{figure}
\centering
\includegraphics[width=\textwidth]{figures/temporal_vs_spatial_effect_size.png}
\caption{Temporal vs spatial effect sizes}
\end{figure}

This result fundamentally reframes urban congestion as a \textbf{temporal synchronization problem}: congestion occurs because everyone travels at the same times, not because certain locations inherently generate more traffic. The LISA hotspots identified above represent temporal bottlenecks---locations where capacity is insufficient during synchronized peak demand---rather than inherently problematic spatial characteristics.

\subsection{Comparative City Analysis}

Cross-city comparison reveals scaling relationships (Figure 6 and Tables 11--12). Congestion does not scale linearly with population: larger cities show both higher absolute congestion and greater peak/off-peak variability, suggesting urban scale amplifies temporal fluctuations.

\begin{figure}
\centering
\includegraphics[width=\textwidth]{figures/boxplot_comparison.png}
\caption{Cross-city comparison of congestion metrics}
\end{figure}

\textbf{Table 11.} Comparative metrics

\begin{longtable}[]{@{}llll@{}}
\toprule\noalign{}
Metric & Jakarta & Bandung & Semarang \\
\midrule\noalign{}
\endhead
\bottomrule\noalign{}
\endlastfoot
Population (million) & 10.5 & 2.5 & 1.8 \\
Mean Jam Factor & 1.599 & 1.537 & 1.299 \\
Off-Peak Period Mean & 1.223 & 1.180 & 1.034 \\
Peak/Off-peak ratio & 1.31x & 1.30x & 1.26x \\
Peak increase (\%) & 30.8\% & 30.2\% & 25.6\% \\
Hotspot density (/km$^2$) & 4.30 & 3.12 & 0.42 \\
\end{longtable}

\textbf{Table 12.} Peak vs Off-Peak Congestion Analysis

\begin{longtable}[]{@{}
  >{\raggedright\arraybackslash}p{(\linewidth - 10\tabcolsep) * \real{0.0909}}
  >{\raggedright\arraybackslash}p{(\linewidth - 10\tabcolsep) * \real{0.1667}}
  >{\raggedright\arraybackslash}p{(\linewidth - 10\tabcolsep) * \real{0.2273}}
  >{\raggedright\arraybackslash}p{(\linewidth - 10\tabcolsep) * \real{0.2273}}
  >{\raggedright\arraybackslash}p{(\linewidth - 10\tabcolsep) * \real{0.1061}}
  >{\raggedright\arraybackslash}p{(\linewidth - 10\tabcolsep) * \real{0.1818}}@{}}
\toprule\noalign{}
City & Peak Mean & Off-Peak Mean & Absolute Diff & Ratio & \% Increase \\
\midrule\noalign{}
\endhead
\bottomrule\noalign{}
\endlastfoot
Jakarta & 1.599 & 1.223 & 0.376 & 1.31x & 30.8\% \\
Bandung & 1.537 & 1.180 & 0.357 & 1.30x & 30.2\% \\
Semarang & 1.299 & 1.034 & 0.265 & 1.26x & 25.6\% \\
\end{longtable}

Jakarta exhibits the largest peak/off-peak differential (30.8\%), reflecting massive daily commuter flows; Bandung shows similar intensification (30.2\%) despite its smaller size, likely due to constrained network geometry; and Semarang the smallest differential (25.6\%), reflecting more distributed logistics-oriented commercial activity.



\section{Discussion}

\subsection{Temporal Patterns and Spatial Clustering}

The dominance of evening peak congestion reflects common Southeast Asian patterns where afternoon activities---school dismissals, commercial operations, and shift changes---coincide with return commutes (Cervero, 2013). The asymmetry between morning and evening peaks (evening $\sim$20\% higher) results from more distributed morning departures, concentrated evening activities, and school/office closing-time overlap. While global Moran's I shows no significant spatial autocorrelation at the road-segment level, this does not preclude local clusters: LISA reveals hotspot patterns geographically coherent with known congested corridors, and congestion propagates through networks consistent with bottleneck queuing theory (Daganzo, 2007). During the evening peak, over half of Jakarta's monitored segments reach moderate congestion (JF > 2.0), while Bandung and Semarang remain predominantly in the light-traffic range. Notably, the H3 scale-sensitivity analysis (Section~\ref{sec:h3_robustness}) reveals that Jakarta does develop weak but statistically significant neighbourhood-scale spatial autocorrelation ($I \approx 0.03$, $p < 0.05$) when congestion is aggregated to hexagonal bins of 174--461\,m diameter---a signal invisible at the individual segment scale due to within-hex noise. This finding suggests that Jakarta's scale and commuter density create modest geographic structure in congestion patterns, a characteristic absent in Bandung and Semarang, consistent with Jakarta's substantially larger network and higher absolute congestion levels.

\subsection{The Failure of Spatial Predictors and Capacity Constraints}

A central and perhaps counterintuitive finding is the \textbf{near-zero correlation} between all spatial characteristics and congestion. Network centrality, land-use intensity (POI density), and road capacity together explain less than 0.02\% of congestion variance. The bottleneck analysis provides perhaps the strongest evidence: three independent tests---aggregate capacity comparison, spatial capacity drop proximity, and local bottleneck gradient---all yield null results. The HERE jam factor's normalization to free-flow speed explains why aggregate capacity effects vanish ($JF \propto 1 - v/v_f$). Crucially, however, the capacity drop analysis tests a \emph{relative} hypothesis (that hierarchy transitions create localized spikes regardless of normalization), and this too fails (p = 0.095--0.830). The combined triangulated evidence from four independent null findings supports the conclusion that spatial capacity constraints do not drive congestion patterns in these cities, consistent with induced demand theory (Duranton \& Turner, 2011). These findings would not have been detectable in prior crowdsourced traffic studies, which typically employed snapshot observations of hundreds to thousands of data points from single cities and relied on descriptive visualization rather than inferential hypothesis testing. The present study's scale---264 million observations across three cities over 11 months---provides the statistical power necessary to confidently reject spatial predictors and quantify the 1,000--4,000$\times$ dominance of temporal over spatial factors.

\subsection{Implications for Traffic Management}

The temporal dominance finding suggests important policy directions, with the caveat that these are based on normalized (relative) congestion metrics and should be confirmed with absolute delay data before informing major infrastructure decisions:

\begin{enumerate}
\item \textbf{Demand management} (staggered work/school hours, flexible work policies, congestion pricing) should be considered alongside infrastructure expansion; distributing departure times across a 2--3 hour window could reduce peak congestion by 30--40\%
\item \textbf{Temporal targeting} of resources to the 16:00--19:00 peak---adaptive signal control, traffic police deployment, incident response---is likely to maximize impact
\item \textbf{Infrastructure decisions} should not rely solely on normalized congestion metrics; betweenness centrality and physical capacity transitions do not predict congestion distribution, suggesting that route diversification and bottleneck removal may have limited effectiveness
\end{enumerate}

\subsection{Limitations}

Key limitations include: (1) HERE API coverage may underrepresent minor roads; (2) the jam factor by design removes absolute capacity effects---the capacity analysis tested relative effects that should still be detectable under normalization, but future studies should use absolute speed data; (3) the 11-month period may not capture long-term trends; (4) traffic segment and OSM edge geometries do not perfectly align, though match rates exceeded 95\%; (5) the ordinal capacity score based on OSM classification is a proxy---lane count, signals, and intersection design are not captured; (6) segment length heterogeneity (200m to >3km contributing equally to spatial statistics) represents a potential MAUP limitation; however, H3 hex-level robustness checks at resolutions 8 and 9 confirm that null POI and centrality correlations are not an artefact of this choice of spatial unit, though they do reveal weak neighbourhood-scale spatial autocorrelation in Jakarta not visible at segment level (Section~\ref{sec:h3_robustness}); and (7) correlational findings do not establish causality.

\subsection{Future Research}

Priority research extensions include: (1) denormalized capacity analysis re-aggregating raw HERE speed data to compute absolute delay ($\Delta t = L \times (1/v - 1/v_f)$) to test whether high-capacity roads carry more total delay despite similar jam factors; (2) temporal capacity interactions examining whether capacity constraints become relevant only during specific peak periods; (3) machine learning approaches combining temporal features with network topology for congestion prediction; and (4) microscopic bottleneck analysis using intersection-level data to test whether signal timing creates localized capacity effects invisible at segment level.



\section{Conclusions}

This study presents a comprehensive spatiotemporal analysis of urban traffic congestion in three Indonesian metropolitan areas using 11 months and 264 million observations. The key findings fundamentally reframe understanding of urban congestion.

\subsection{Primary Finding: Temporal Dominance}

\textbf{Relative congestion is fundamentally a temporal phenomenon.} Time-of-day explains 15--24\% of jam factor variance ($\eta^2$ from ANOVA), while all spatial predictors combined explain less than 0.02\%---a 1,000--4,000$\times$ difference. An important caveat is that the jam factor captures demand-to-capacity ratio (normalized metric); whether temporal dominance also holds for absolute delay remains an open question requiring raw speed and volume data.

\subsection{Secondary Findings}

\begin{enumerate}
\item \textbf{Evening peak dominance:} Congestion during 16:00--19:00 exceeds daily averages by $\sim$40\%, with Jakarta reaching a mean jam factor of 2.01
\item \textbf{Local clustering despite global randomness:} Global Moran's I is non-significant at the road-segment level across all three cities, but LISA identifies local hotspots ($\sim$10\% of segments) representing peak-hour capacity bottlenecks; H3 hex-level analysis additionally reveals weak but consistent neighbourhood-scale spatial autocorrelation in Jakarta ($I = 0.030$, $p = 0.033$ at resolution~8, $\approx$461\,m hex diameter), a signal absent in Bandung and Semarang
\item \textbf{Scale effects:} Larger cities show higher absolute congestion and greater peak/off-peak differentials (Jakarta 30.8\%, Semarang 25.6\%)
\item \textbf{Network topology irrelevance:} Betweenness centrality correlates near-zero with congestion (r < 0.02)
\item \textbf{Capacity constraint irrelevance:} Neither aggregate road capacity ($|d|$ = 0.03--0.07, p > 0.05), spatial capacity drops (p = 0.10--0.83), nor local bottleneck gradients ($|d|$ < 0.05) predict congestion levels
\end{enumerate}

\subsection{Policy Implications and Contribution}

These findings support demand management strategies (staggered hours, flexible work, congestion pricing) alongside infrastructure expansion, given the strong temporal dominance in relative congestion. Temporal targeting of resources to the 16:00--19:00 peak period is likely to maximize impact regardless of metric choice. Infrastructure investment decisions should incorporate absolute delay analysis rather than relying solely on normalized metrics.

This research contributes empirical evidence demonstrating that urban congestion is primarily a temporal coordination failure rather than a spatial infrastructure or capacity constraint problem. A parallel software contribution is an open-source Python package that provides a fully programmable, end-to-end traffic analysis pipeline: Python-native HERE API collection supersedes the R-based approach used for this study's data, and the remaining stages---large-scale aggregation, geostatistical and network analysis, HPC-enabled centrality computation, and automated manuscript generation---run entirely within Python. This pipeline is directly transferable to other cities and commercial traffic data sources. The novel graph-based capacity drop detection framework provides a replicable method for testing the bottleneck hypothesis in other urban contexts, and the finding that time explains 1,000$\times$ more variance than space---corroborated by four independent null capacity tests and confirmed robust to spatial scale by Uber H3 hexagonal aggregation---has significant implications for transportation policy worldwide.



\section*{Acknowledgments}

Computational resources were provided by Pusdatin Bappenas (The Data and Information Center, The Ministry of National Development Planning), which granted access to NVIDIA H100 GPUs on their High Performance Computing cluster.



\section*{Declaration of Competing Interests}

The authors declare no competing financial interests or personal relationships that could have influenced this work.



\section*{Author Contributions}

\textbf{Firman Hadi:} Conceptualization, Methodology, Software, Data curation, Formal analysis, Writing -- original draft. \textbf{Yasser Wahyuddin:} Methodology, Formal analysis, Writing -- original draft. \textbf{L.M. Sabri:} Formal analysis, Writing -- review \& editing. \textbf{Agung Indrajit:} Writing -- review \& editing.



\section*{Declaration of Generative AI and AI-assisted Technologies in the Writing Process}

During the preparation of this work, the authors used Claude (Anthropic) for data processing pipeline development, evaluation script implementation, and statistical analysis automation. The authors reviewed and edited all outputs and take full responsibility for the content of this publication.



\section*{Ethics Statement}

This research used publicly available satellite imagery and did not involve human participants, animal subjects, or sensitive personal data. No ethical approval was required.



\section*{Funding}

This research did not receive any specific grant from funding agencies in the public, commercial, or not-for-profit sectors.



\section*{Data Availability}

Traffic flow data were obtained from the HERE Traffic API. Aggregated datasets and analysis code are available at: https://github.com/firmanhadi21/traffic-analyses

Street network data are publicly available from OpenStreetMap (www.openstreetmap.org).



\section*{References}

Alkaabi, K.,Raza, M., Qassemi, E., Alderei, H., Alderei, M \& Almheiri, S. (2024). Using crowd-sourced traffic data and open-source tools for urban congestion analysis. \emph{Case Studies on Transport Policy}, 24, 101261.

Anderson, T. K. (2009). Kernel density estimation and K-means clustering to profile road accident hotspots. \emph{Accident Analysis \& Prevention}, 41(3), 359-364.

Anselin, L. (1995). Local indicators of spatial association---LISA. \emph{Geographical Analysis}, 27(2), 93-115.

Arnott, R., de Palma, A., \& Lindsey, R. (1993). A structural model of peak-period congestion: A traffic bottleneck with elastic demand. \emph{American Economic Review}, 83(1), 161-179.

Asian Development Bank. (2019). \emph{Asian Development Outlook 2019: Strengthening Disaster Resilience}. Manila: ADB.

Barrington-Leigh, C., \& Millard-Ball, A. (2020). Global trends toward urban street-network sprawl. \emph{Proceedings of the National Academy of Sciences}, 117(4), 1941-1950.

Batty, M. (2013). \emph{The New Science of Cities}. Cambridge, MA: MIT Press.

Gedamu, W.T., Plank-Wiedenbeck, U., \& Wodajo, B.T. (2024). A spatial autocorrelation analysis of road traffic crash by severity using Moran's I spatial statistics: A comparative study of Addis Ababa and Berlin cities. \emph{Accident Analysis \& Prevention}, 197, 107458.

Benjamini, Y., \& Hochberg, Y. (1995). Controlling the false discovery rate: A practical and powerful approach to multiple testing. \emph{Journal of the Royal Statistical Society: Series B}, 57(1), 289-300.

Boeing, G. (2017). OSMnx: New methods for acquiring, constructing, analyzing, and visualizing complex street networks. \emph{Computers, Environment and Urban Systems}, 65, 126-139.

Boeing, G. (2018). Urban street network centrality. \emph{Geoff Boeing Blog}. Retrieved from geoffboeing.com/2018/01/urban-street-network-centrality/

Boeing, G. (2022). Street network models and indicators for every urban area in the world. \emph{Geographical Analysis}, 54(3), 519-535.

Boeing, G. (2025). Modeling and analyzing urban networks and amenities with OSMnx. \emph{Geographical Analysis}, 57(1), 1-18.

BPS Bandung. (2015). \emph{Bandung City in Figures 2015}. Bandung: Badan Pusat Statistik Kota Bandung.

Cai, X., Xiao, Y., \& Zhang, L. (2024). Traffic congestion cause identification method for urban main roads. \emph{Transportation Research Record}, 2678(8), 456-470.

Castrol. (2015). \emph{Castrol Magnatec Stop-Start Index}. London: BP.

Cervero, R. (2013). Linking urban transport and land use in developing countries. \emph{Journal of Transport and Land Use}, 6(1), 7-24.

Cui, H., Dong, X., Su, Y., Zhu, M., \& Ma, X. (2024). Exploring spatiotemporal patterns of frequently congested urban road segments based on multi-source data: A case study of China's super-large cities. \emph{Transportation Research Record}, 2678(8), 789-804.

Daganzo, C. F. (2007). Urban gridlock: Macroscopic modeling and mitigation approaches. \emph{Transportation Research Part B}, 41(1), 49-62.

Dasgupta, S., Wheeler, D., Khaliquzzaman, M., \& Huq, M. (2021). Siting Priorities for Congestion-Reducing Projects in Dhaka.

Duranton, G., \& Turner, M. A. (2011). The fundamental law of road congestion: Evidence from US cities. \emph{American Economic Review}, 101(6), 2616-2652.

Ermagun, A., \& Levinson, D. (2018). Spatiotemporal traffic forecasting: Review and proposed directions. \emph{Transport Reviews}, 38(6), 786-814.

Gakenheimer, R. (1999). Urban mobility in the developing world. \emph{Transportation Research Part A}, 33(7-8), 671-689.

Getis, A., \& Ord, J. K. (1992). The analysis of spatial association by use of distance statistics. \emph{Geographical Analysis}, 24(3), 189-206.

Gao, S., Wang, Y., Gao, Y., \& Liu, Y. (2013). Understanding urban traffic-flow characteristics: A rethinking of betweenness centrality. \emph{Environment and Planning B}, 40(1), 135-153.

Dasgupta, S., Lall, S., \& Wheeler, D. (2021). Spatiotemporal analysis of traffic congestion, air pollution, and exposure vulnerability in Tanzania. \emph{Science of the Total Environment}, 778, 147114.

HERE Technologies. (2024). \emph{HERE Traffic Analytics: Speed Data Enhancements}. Retrieved from here.com/learn/blog/here-traffic-analytics-speed-data-enhancements.

Jakarta Environmental Agency. (2022). \emph{Smart Traffic Management System Implementation Report 2022}. Jakarta: Dinas Lingkungan Hidup DKI Jakarta.

Jenelius, E., \& Koutsopoulos, H. N. (2015). Probe vehicle data sampled by time or space: Consistent travel time allocation and estimation. \emph{Transportation Research Part B}, 71, 120-137.

%    Jiang, B., \& Claramunt, C. (2004). Topological analysis of urban street networks. \emph{Environment and Planning B}, 31(1), 151-162.

HERE Technologies. (2023). \emph{HERE Traffic API Developer Guide}. Retrieved from developer.here.com.

Joewono, T. B., \& Kubota, H. (2008). Paratransit service in Indonesia: User satisfaction and future choice. \emph{Transportation Planning and Technology}, 31(3), 325-345.

Kirkley, A., Barbosa, H., Barthelemy, M., \& Ghoshal, G. (2018). From the betweenness centrality in street networks to structural invariants in random planar graphs. \emph{Nature Communications}, 9, 2501.

Kutzbach, M. J. (2009). Motorization in developing countries: Causes, consequences, and effectiveness of policy options. \emph{Journal of Urban Economics}, 65(2), 154-166.

%Louail, T., Lenormand, M., Cantu Ros, O. G., Picornell, M., Herranz, R., Frias-Martinez, E., Ramasco, J. J., \& Barthelemy, M. (2015). From mobile phone data to the spatial structure of cities. \emph{Scientific Reports}, 4, 5276.

%Louf, R., \& Barthelemy, M. (2014). A typology of street patterns. \emph{Journal of the Royal Society Interface}, 11(101), 20140924.

Leduc, G. (2008). Road traffic data: Collection methods and applications. \emph{JRC Technical Notes}, European Commission.

%Marshall, W. E., Piatkowski, D. P., \& Garrick, N. W. (2018). Community design, street networks, and public health. \emph{Journal of Transport \& Health}, 1(4), 326-340.

Li, X., Cui, J., An, S., \& Parsafard, M. (2018). Stop-and-go traffic analysis: Theoretical properties, environmental impacts and oscillation mitigation. \emph{Transportation Research Part B}, 70, 319-339.

Ma, Y., Zhang, H., \& Liu, W. (2024). Traffic state prediction method based on spatial--temporal data mining of floating car data. \emph{Computer-Aided Civil and Infrastructure Engineering}, 39(5), 712-730.

Moudon, A. V., Lee, C., Cheadle, A. D., Collier, C. W., Johnson, D., Schmid, T. L., \& Weather, R. D. (2006). Operational definitions of walkable neighborhood: Theoretical and empirical insights. \emph{Journal of Physical Activity and Health}, 3(s1), S99-S117.

Moran, P. A. P. (1950). Notes on continuous stochastic phenomena. \emph{Biometrika}, 37(1-2), 17-23.

Openshaw, S. (1984). \emph{The Modifiable Areal Unit Problem}. Concepts and Techniques in Modern Geography, 38. Norwich: Geo Books.

Ord, J. K., \& Getis, A. (1995). Local spatial autocorrelation statistics: Distributional issues and an application. \emph{Geographical Analysis}, 27(4), 286-306.

Porta, S., Crucitti, P., \& Latora, V. (2006). The network analysis of urban streets: A primal approach. \emph{Environment and Planning B}, 33(5), 705-725.

Pojani, D., \& Stead, D. (2015). Sustainable urban transport in the developing world: Beyond megacities. \emph{Sustainability}, 7(6), 7784-7805.

Rey, S. J., \& Anselin, L. (2010). PySAL: A Python library of spatial analytical methods. In M. M. Fischer \& A. Getis (Eds.), \emph{Handbook of Applied Spatial Analysis} (pp.~175-193). Berlin: Springer.

Rempe, F., Huber, G., \& Bogenberger, K. (2016). Spatio-temporal congestion patterns in urban traffic networks. \emph{Transportation Research Procedia}, 15, 513-524.

Saberi, M., Hamedmoghadam, H., Ashfaq, M., Hosseini, S. A., Gu, Z., Shafiei, S., Nair, D. J., Dixit, V., Gardner, L., Waller, S. T., \& Gonz\'{a}lez, M. C. (2020). A simple contagion process describes spreading of traffic jams in urban networks. \emph{Nature Communications}, 11, 1616.

Song, J.,  Zhoo, C., Zhong, S., Nielsen,  T.A.S.,\& Prishchepov, A.V. (2019). Mapping spatio-temporal patterns and detecting the factors of traffic congestion with multi-source data fusion and mining techniques. \emph{Computers, Environment and Urban Systems}, 77, 101364.

Susilo, Y. O., Joewono, T. B., Santosa, W., \& Parikesit, D. (2007). A reflection of motorization and public transport in Jakarta metropolitan area. \emph{IATSS Research}, 31(1), 59-68.

Syaiful, S., \& Fadlilah, N. (2020). Exploring the effectiveness of demand management policy in reducing traffic congestion and environmental pollution: Car-free day and odd-even plate measures for Bandung city in Indonesia. \emph{Case Studies on Transport Policy}, 8(3), 713-720.

Tobler, W. R. (1970). A computer movie simulating urban growth in the Detroit region. \emph{Economic Geography}, 46(sup1), 234-240.

TomTom. (2023). \emph{TomTom Traffic Index 2023}. Amsterdam: TomTom International.

TomTom. (2025). \emph{Traffic Stats API Documentation}. Retrieved from developer.tomtom.com/traffic-stats.

Wang, Y., Gu, Y., Dou, M., \& Qiao, M. (2016). Using spatial semantics and interactions to identify urban functional regions. \emph{ISPRS International Journal of Geo-Information}, 7(4), 130.


WHO. (2021). \emph{WHO Global Air Quality Guidelines}. Geneva: World Health Organization.

World Bank. (2019). \emph{Bandung Urban Transport Diagnostic}. Washington, DC: World Bank Group.

World Bank. (2021). \emph{Jakarta Urban Mobility Assessment}. Washington, DC: World Bank Group.

World Bank. (2022). \emph{Indonesia Mass Transit Project}. Washington, DC: World Bank Group.

Yuan, Y., Zhang, X., \& Li, S. (2025). Review and prospect of floating car data research in transportation. \emph{Communications in Transportation Research}, 6, 100117.

Zhang, L., Wang, J., \& Liu, X. (2022). Spatiotemporal analysis of urban road congestion during and post COVID-19 pandemic in Shanghai, China. \emph{Scientific Reports}, 12, 2024.

Zhong, C., Arisona, S. M., Huang, X., Batty, M., \& Schmitt, G. (2014). Detecting the dynamics of urban structure through spatial network analysis. \emph{International Journal of Geographical Information Science}, 28(11), 2178-2199.



\section*{Supplementary Materials}

Additional figures, detailed statistical outputs, and analysis code are available at: https://github.com/firmanhadi21/traffic-analyses

\end{document}
